\secciongrande{Lenguajes de programación en entorno cliente dos entornos de ejecución}

\resumen{
  La programación web requiere diferenciar entre entornos cliente y servidor. El entorno cliente se ejecuta en el navegador del usuario, mientras que el servidor procesa peticiones y gestiona bases de datos. La elección del lenguaje depende de las ventajas e inconvenientes, siendo crucial la complementariedad entre ambos entornos para sitios dinámicos.
}

\definicion{Entorno cliente}{Entorno de ejecución donde el código es procesado por el dispositivo del usuario, típicamente un navegador web.}

\definicion{Entorno servidor}{Entorno de ejecución donde el código reside y se procesa en el servidor remoto, gestionando lógica de negocio y bases de datos.}

\subtituloA{Diferenciación y elección de lenguajes}

\lista{
  \textbf{Entorno cliente:} Ejecutado por el cliente en el navegador web, enfocado en la interactividad del usuario. |
  \textbf{Entorno servidor:} Ejecutado en el servidor, responsable del almacenamiento de datos y la lógica de negocio compleja.
}

\nota{Para el examen}{
Es vital que el desarrollador evalúe las ventajas e inconvenientes de cada lenguaje, ya que la decisión final es contextual y depende de su pericia profesional.
}
